\section{Experimental Setup}
\label{sec:setup}

\subsection{Training configuration}
\label{sec:setup:training}

Table~\ref{tab:hparams} lists the optimisation settings, which are identical
across the three Transformer variants.

\begin{table}[htbp]
\centering
\caption{Training configuration for the Transformer variants.}
\label{tab:hparams}
\begin{tabular}{ll}
\toprule
Setting & Value \\
\midrule
Optimiser              & AdamW~\cite{Loshchilov2019AdamW} \\
Learning rate          & $10^{-4}$, cosine decay to $5\times10^{-7}$ \\
Warm-up                & 10 epochs \\
Weight decay           & 0.05 (excluding norms, biases, token parameters) \\
Gradient clipping      & 1.0 \\
Batch size             & 4, with 4 accumulation steps (effective 16) \\
Epoch budget           & 100 (early stopped) \\
Precision              & bf16 mixed precision \\
Memory                 & gradient checkpointing \\
Variable weights       & uniform $[1,1,1,1,1]$ \\
Model selection        & validation MAE, patience 40 \\
Early stopping         & no improvement for 20 validation checks \\
Random seed            & 42 \\
\bottomrule
\end{tabular}
\end{table}

Weight decay is assigned by module type and parameter role rather than by
name matching: normalisation parameters, biases and token-like parameters
(mask tokens, learned state embeddings) are excluded.

\subsection{Model variants}
\label{sec:setup:variants}

Four Transformer variants are compared, plus a recurrent baseline and two
non-learned baselines. The \textbf{MAE Transformer} is the reference
configuration: encoder spatial attention enabled, trained at mask ratio~0.5.
The \textbf{Probabilistic MAE Transformer} changes only the loss, from Huber
to a heteroscedastic Gaussian negative log-likelihood with a predicted
variance head, isolating the effect of the training objective. The
\textbf{Dense Transformer} keeps spatial attention but trains at mask
ratio~0 instead of 0.5, isolating the training mask ratio. The
\textbf{Spatially Blind Transformer} removes the encoder's spatial sub-layer
\emph{and} makes the decoder station-local, so that every station's
prediction depends only on its own encoder tokens --- the fully
station-independent control against which spatial modelling is judged. The
\textbf{LSTM Baseline} provides a non-Transformer, purely recurrent
reference at the same task. \textbf{Last-Value Persistence} and \textbf{24h
Persistence} are non-learned baselines, described further in
Section~\ref{sec:setup:metrics}, and are included in Table~\ref{tab:scenarios}
for completeness.

Table~\ref{tab:scenarios} lists the characteristics of every scenario. The
MAE Transformer and Spatially Blind Transformer differ in encoder spatial
attention \emph{and} decoder scope simultaneously, so together they isolate
the total contribution of cross-station information, not one architectural
component in isolation; the MAE Transformer and Dense Transformer differ
only in training mask ratio, isolating that variable on its own.

\begin{table}[htbp]
\centering
\caption{Characteristics of the evaluated scenarios. Parameter counts and
best epoch are read directly from the checkpoints; the persistence baselines
involve no training.}
\label{tab:scenarios}
\begin{tabular}{lccccrr}
\toprule
Scenario & Spatial attn. & Decoder scope & Train MR & Loss & Params & Best epoch \\
\midrule
MAE Transformer               & yes & cross-station    & 0.5 & Huber        & 25.63~M & 40 \\
Probabilistic MAE Transformer & yes & cross-station    & 0.5 & Gaussian NLL & 25.63~M & 34 \\
Dense Transformer              & yes & cross-station    & 0.0 & Huber        & 25.63~M & 47 \\
Spatially Blind Transformer    & no  & station-local    & 0.0 & Huber        & 20.89~M & 26 \\
LSTM Baseline                  & n/a & n/a (recurrent)  & n/a & Huber        & 21.11~M & 40 \\
Last-Value Persistence         & --  & --               & --  & --           & --      & -- \\
24h Persistence                 & --  & --               & --  & --           & --      & -- \\
\bottomrule
\end{tabular}
\end{table}

The 4.74~M parameter difference between the MAE Transformer and the
Spatially Blind Transformer corresponds exactly to eight spatial
multi-head-attention sub-layers plus their layer normalisations, confirming
that the ablation removed the intended component and nothing else; the
station-local decoder re-wires attention scope rather than adding
parameters, so it does not show up in this count.

\paragraph{Ablation scope.}
\label{sec:method:ablations}

Unlike an ablation that disables only the encoder's spatial sub-layer while
leaving the decoder's cross-attention over all stations intact --- which
would still let information flow between stations through the decoder ---
the Spatially Blind Transformer additionally sets
\texttt{station\_local\_decoder}, closing that pathway too. Its comparison
against the MAE Transformer therefore measures the total effect of
cross-station information, both in the encoder and in the decoder, rather
than the encoder's spatial sub-layer alone. A configuration that disables
only the encoder's spatial attention while keeping the decoder global was
not trained for this project and is noted as a direction for future work in
Section~\ref{sec:setup:known-issues}.

\subsection{Data splits and window sampling}
\label{sec:setup:splits}

Splits are chronological: training 2017--2021, validation 2022, test
2023--2024. No future observation informs a past prediction, and normalisation
statistics are computed on the training years alone
(Section~\ref{sec:data:preparation}).

Training draws windows at random from the pool of valid start times, with an
epoch defined as 10{,}000 sampled windows rather than a full pass. Because the
station mask is redrawn on every forward pass, revisiting the same time window
in a later epoch presents a different task, so resampling is not redundant.
Validation uses non-overlapping blocks over 2022. Testing uses sliding windows
at a 90-minute stride over 2023--2024, giving 11{,}684 windows.

\subsection{Metrics}
\label{sec:setup:metrics}

Errors are reported as MAE, RMSE, bias and $R^2$, in two unit systems.
Normalised units allow comparison across variables of different physical scale.
Physical units are obtained by applying the originating station's own standard
deviation,
%
\begin{equation}
  \varepsilon^{\text{phys}}_{s,v}
    = \big(\hat{y}_{s,v} - y_{s,v}\big)\,\sigma_{s,v},
  \label{eq:denorm}
\end{equation}
%
where the station mean cancels in the difference. Using a network-averaged
standard deviation instead biases variables whose spread varies across stations;
for wind this spread is a factor of 14, so the choice is not cosmetic.

Skill is reported against persistence --- the last observed value of the input
window, carried forward to every lead time~\cite{Murphy1988SkillScores}. Wind
speed and direction are reconstructed from the predicted $u$/$v$ components
after per-station denormalisation, with direction error reported only for wind
speeds above 3~m\,s$^{-1}$, since direction is undefined in calm air.

\subsection{Evaluation protocol}
\label{sec:setup:protocol}

Each model is evaluated at two mask ratios. At $\text{MR} = 0$ all stations are
visible and the task is pure forecasting; this is the only regime the LSTM
supports, and therefore the only setting in which all four models are
comparable. At $\text{MR} = 0.5$ half the stations are hidden, and predictions
for hidden and visible stations are reported separately, because they are
different tasks: a hidden station must be reconstructed from its neighbours,
whereas a visible station has its own history available.

Results are broken down by lead time (0 to 6~h in 30-minute steps), by variable,
and by station. The $\Delta = 0$ column is reported separately from $\Delta > 0$
throughout, since at $\Delta = 0$ a visible station can copy its own last
observation and the persistence residual makes this explicit; pooling it with
genuine forecast horizons inflates apparent accuracy.

\subsection{Reproducibility}
\label{sec:setup:repro}

Training seeds Python, NumPy and PyTorch from a single value (default 42).
Full determinism is not enforced (\texttt{deterministic=False}), so exact
bitwise reproduction of a training run is not guaranteed on GPU.

The full configuration of every run, including the normalisation statistics, is
serialised into its checkpoint, so any trained model can be rebuilt exactly from
its own file. All prediction dumps share bit-identical targets, masks and time
axes (verified by assertion in the data-audit notebook), so model comparisons
are paired at the observation level. A test suite of 17 modules covers the leak
guard, shape contracts, embedding numerics, configuration round-tripping and
the per-station denormalisation.

\subsection{Known issues affecting interpretation}
\label{sec:setup:known-issues}

Three implementation details are recorded here because they qualify the
results rather than the method.

\emph{Train/evaluation mask mismatch.} The MAE Transformer and the
Probabilistic MAE Transformer train at $\text{MR} = 0.5$ but are
model-selected on validation with all stations visible ($\text{MR} = 0$).
The encoder therefore sees roughly 1{,}872 tokens during training and 3{,}720
during validation --- a token count it never receives a gradient update on.
The Dense Transformer exists in part to isolate this effect, since it trains
and validates at the same mask ratio.

\emph{Evaluation-mask seeding.} The station mask drawn at $\text{MR} = 0.5$
evaluation time is sampled from the global random-number generator on every
forward pass. Earlier prediction dumps (before 2026-08-25) drew this mask
without a fixed seed, so two evaluation runs of different models masked
different station subsets and cross-model masked-station comparisons were
unpaired. \texttt{test.py} now re-seeds this generator once per mask ratio
before evaluation, so masked-station comparisons across scenarios are paired
whenever the same seed, batch size, index mode and stride are used. Results
in this report use dumps produced after this fix; any figure computed from a
dump predating it should be regenerated before being quoted.

\emph{No partial spatial-attention ablation.} As noted in
Section~\ref{sec:method:ablations}, a configuration that disables only the
encoder's spatial sub-layer while keeping the decoder's cross-station
attention intact was not trained for this project. The Spatially Blind
Transformer closes both pathways at once, so it cannot separate "the encoder
no longer mixes stations" from "the decoder no longer mixes stations" --- it
answers only their combined effect.
